% DOCUMENT CLASS %
\documentclass[a4paper, 12pt]{extarticle}   % for A4 Paper, 12pt fontsize and article spacing
\usepackage[left=1.5cm,right=1.5cm,top=2cm,bottom=2cm]{geometry}    % Make the page a little bigger

% PACKAGES%
\usepackage{amsfonts}   % For a basic mathfont (many will be replaced by stix)
\usepackage{amsmath}    % For basic math symbols
\usepackage{bbm}        % For \mathbbm (better version of \mathbb)
\usepackage{mathrsfs}   % For \mathscr
\usepackage{hyperref}   % For footnotes
\usepackage{csquotes}   % For \textquote{}
\usepackage{IEEEtrantools}  % For better alignments
\usepackage{tikz}   % For a macro
\usetikzlibrary{arrows.meta, positioning, calc}	% For protocol diagrams
\usepackage{listings} % For code
\usepackage{xcolor}   % For ... code
\usepackage{microtype}  % For better paragraph formatting
\usepackage{amsthm}   % For custom environments

% FONT CONFIGURATION %
\usepackage[T1]{fontenc}
\usepackage{palatino}

% BASIC SETS %
\newcommand*{\R}{\mathbbm R}    % Set of real numbers
\newcommand*{\N}{\mathbbm N}    % Set of natural numbers (beginning at 0)
\newcommand*{\Z}{\mathbbm Z}    % Set of integers
\newcommand*{\Q}{\mathbbm Q}    % Set of rational numbers

% MACROS %
\newcommand*{\qedbox}{\null\nobreak\hfill\ensuremath{\square}} % Macro for a QED-Box
% for a funny ¯\_(ツ)_/¯-Emoji
\newcommand{\shrug}[1][]{%
    \begin{tikzpicture}[baseline,x=0.8\ht\strutbox,y=0.8\ht\strutbox,line width=0.125ex,#1]
    \def\arm{(-2.5,0.95) to (-2,0.95) (-1.9,1) to (-1.5,0) (-1.35,0) to (-0.8,0)};
    \draw \arm;
    \draw[xscale=-1] \arm;
    \def\headpart{(0.6,0) arc[start angle=-40, end angle=40,x radius=0.6,y radius=0.8]};
    \draw \headpart;
    \draw[xscale=-1] \headpart;
    \def\eye{(-0.075,0.15) .. controls (0.02,0) .. (0.075,-0.15)};
    \draw[shift={(-0.3,0.8)}] \eye;
    \draw[shift={(0,0.85)}] \eye;
    % draw mouth
    \draw (-0.1,0.2) to [out=15,in=-100] (0.4,0.95); 
    \end{tikzpicture}
}

% BASIC CONFIGURATION %
\pagestyle{plain}
\allowdisplaybreaks

% CODE CONFIGURATIONS %
\definecolor{codepurple}{HTML}{7f0055}
\definecolor{comment}{RGB}{0,128,0}
\definecolor{number}{RGB}{9,136,90}
\definecolor{string}{RGB}{163,21,21}

\lstset{
    language=python,
    emph={*, do, od, for, if, fi, then, else, to, def, in, forall, exists},
    firstnumber=1,
    basicstyle=\small\ttfamily,
    breaklines=true,
    breakatwhitespace=false,
    frame=tb,
    captionpos=b,
    %float=t,
    %floatplacement=t,
    abovecaptionskip=0.5cm,
    numbers=left,
    numberstyle=\tiny,
    keywordstyle=\bfseries\color{codepurple},
    emphstyle=\bfseries\color{codepurple},
    mathescape=true,
    keepspaces=true,
    showspaces=false,
    showstringspaces=false,
    showtabs=false,
    stringstyle=\color{string},
    commentstyle=\itshape\color{comment},
    upquote=true,
    texcl=true
}

% EXERCISE ENVIRONMENT %
\theoremstyle{definition}
\newtheorem{exercise}{Exercise}[section]
\setcounter{section}{1} % Don't use \section and this works naturally

% TITLE CONFIGURATION %
\newcommand{\ex}{0}

\title{Data Visualization \\ \small Exercise Sheet \ex}
\date{Winter Term 2025/26}
\author{
    Henri Heyden \\%
    \small stu240825%
    \and
    Julia Schröder \\%
    \small stu252018%
    \and
    Rebecca Ziebell \\%
    \small stu250567%
}

\renewcommand{\ex}{9}
\setcounter{section}{\ex}

\renewcommand{\labelenumi}{\alph{enumi})}

\usetikzlibrary{graphs}
\usetikzlibrary{graphdrawing}
\usegdlibrary{circular}

\begin{document}
    \maketitle
    
    \begin{exercise}[Network Visualization] \hfill
        \begin{enumerate}
            \item \hfill
                \begin{center}
                    \tikz \graph [simple necklace layout, node distance=3cm] {
                        A -- B,
                        A -- C,
                        B -- C,
                        B -- D,
                        C -- E,
                        D -- E
                    };
                \end{center}
            \item A good criterion is minimized edge crossing---at best the graph has a planar embedding.
                  In circular layout, we can permutate the ordering of the vertices to achieve minimized edge crossing.
                  Additionally, the edges in the graph can be bend to improve on this even further.
            \item We can bend some edges to the right and some edges to the left:
                \begin{center}
                    \tikz \graph [node distance=3cm] {
                        A,
                        B,
                        C,
                        D,
                        E,
                        A -- [bend left] B,
                        A -- [bend right] C,
                        B -- [bend left] C,
                        B -- [bend right] D,
                        C -- [bend left] E,
                        D -- [bend right] E
                    };
                \end{center}
            \item Edge bundling clusters similar edges together into shared paths to reduce visual clutter and reveal high-level connection patterns.
        \end{enumerate}
    \end{exercise}
\end{document}
